\begin{table}[t]
\centering
\caption{Women's share among classified sightings in collected street imagery by city.}
\label{tab:city_summary}
\scriptsize
\begin{threeparttable}
\begin{tabular}{@{\extracolsep{0pt}}lrrrr}
\toprule
 & Mumbai & Navi Mumbai & Bangalore & Delhi \\
\midrule
Images annotated & 2,809 & 1,999 & 2,015 & 2,417 \\
Collection days (clusters) & 11 & 9 & 13 & 14 \\
\quad effective (Kish) & 6.2 & 5.6 & 9.1 & 12.4 \\
Adult pedestrian sightings & 12,331 & 4,645 & 2,311 & 3,020 \\
\addlinespace
Pedestrians: women, person-weighted (\%) & 19.1 [17.4, 20.8] & 18.0 [15.7, 20.4] & 27.4 [25.7, 29.1] & 19.8 [17.3, 22.2] \\
Pedestrians: women, image-level mean (\%) & 18.3 [16.4, 20.2] & 18.7 [16.4, 21.0] & 28.4 [25.7, 31.2] & 19.4 [16.7, 22.1] \\
\addlinespace
Women per 1,000 men & 236 & 220 & 377 & 246 \\
\addlinespace
All modes: women, person-weighted (\%) & 15.8 [14.8, 16.9] & 13.6 [12.4, 14.9] & 16.0 [14.6, 17.5] & 12.0 [10.5, 13.5] \\
Two-wheelers: women (\%) & 8.2 & 5.3 & 8.8 & 5.9 \\
All modes: women, image-level mean (\%) & 14.5 [13.2, 15.7] & 12.9 [11.6, 14.3] & 14.5 [12.1, 17.0] & 11.2 [9.4, 12.9] \\
\bottomrule
\end{tabular}
\begin{tablenotes}[flushleft]
\item The first two share rows and the sex ratio describe pedestrians only, which is the primary estimand: the share of classified pedestrian sightings coded as women. The all-mode rows additionally include two-wheeler riders.
\item 95\% confidence intervals use collection-day cluster-robust standard errors and quantify variability across observed fieldwork days; the route imagery is not a population sample.
\item Collection days are the clustering unit, but they contributed very unequal numbers of frames, so the effective count is well below the raw count. The Kish effective number, $(\sum n_c)^2 / \sum n_c^2$, is the one that governs interval coverage. In Mumbai and Navi Mumbai it falls to about six and five, and intervals for those cities should be read as indicative.
\item Person-weighted estimates weight each classified sighting equally; image-level means weight each frame equally.
\item Sex ratio is women per 1,000 men. Gender inferred from visible appearance.
\item ``10+'' is represented by its known minimum, 11; Table~\ref{tab:topcode} reports sensitivity.
\end{tablenotes}
\end{threeparttable}
\end{table}
